\chapter{El \emph{streaming backend}: \emph{canonical store}}
\label{cap:streaming-mirror}

\begin{intuicion}
\texttt{add\_relation} hace un \emph{single write} al
\texttt{streaming} backend ---el almacenamiento canónico desde
stage 2c part 2 (capítulo~\ref{cap:streaming-canonical}). El
backend es una estructura \emph{degree-adaptive per-vértice}
(\texttt{Inline} para grado bajo, \texttt{Vec<StreamEdge>} para
grado alto, \texttt{Mmap} para vértices spilled a disco) con
\emph{auto-spill incremental} cuando el budget global cruza un
umbral. Lecturas son \emph{zero-copy}: el matcher hot-path toma
\texttt{neighbors\_arc(sid)} + \texttt{guard.read()} +
\texttt{as\_slice()} y consume un \texttt{\&[StreamEdge]} con la
misma forma que el legacy \texttt{\&[CompactRelation]}. Este
capítulo explica el shape degree-adaptive y la mecánica del
auto-spill; el capítulo siguiente analiza cómo se hizo el migrate
del matcher para que la antigua copia paralela
(\texttt{SpillableEdgeStore}) pudiera eliminarse.
\end{intuicion}

\section{Por qué un \emph{streaming backend} y no un \texttt{Vec}}

La forma natural de un edge store sería un \texttt{Vec<CompactRelation>}
sortable + binary-searchable. Eso fue el \texttt{SpillableEdgeStore}
legacy. Tiene tres limitaciones para los consumers actuales:

\begin{itemize}
    \item \textbf{Sólo binary-searchable post-\texttt{finalize}}:
    durante ingesta el slice no está sorted, así que un matcher
    incremental que querría re-evaluar una hipótesis con cada
    arista nueva no puede usar \texttt{partition\_point} sobre el
    rango temporal.
    \item \textbf{No expone \emph{per-vertex} access barato}: cada
    \texttt{get\_edges(sid)} hace un \emph{HashMap lookup}; los
    consumers que iteran muchos vértices de un grafo (NLF build,
    khop\_reach build) pagan ese lookup por vértice.
    \item \textbf{Spill-to-disk one-shot}: spillea cuando el buffer
    cruza un umbral, pero post-spill no acepta más writes hasta el
    próximo finalize. No sirve para \emph{ingesta continua}.
\end{itemize}

El \emph{streaming backend} resuelve los tres simultáneamente:
per-vertex \emph{slot} en un \texttt{HashMap<StrId, Arc<RwLock<NeighborList>>>},
\emph{auto-spill incremental} con archivos round-numbered, y un
\emph{per-vertex slice} sorted-by-timestamp post-\texttt{finalize}
que el matcher consume con \texttt{partition\_point} igual que el
viejo store.

\section{La estructura: \texttt{NeighborList} degree-adaptive}

Cada vértice del streaming backend tiene una \texttt{NeighborList}
con tres variantes:

\begin{lstlisting}[language=Rust]
pub enum NeighborList {
    Inline(SmallVec<[StreamEdge; 8]>),    // <= 8 aristas
    Vec(Vec<StreamEdge>),                  // > 8 aristas, in-RAM
    Mmap { handle: Arc<MmapHandle>,        // spilled
           byte_offset: usize, len: usize },
}
\end{lstlisting}

\noindent
La elección depende del estado del vértice:

\begin{itemize}
    \item \textbf{Inline}: 8 aristas \emph{en la propia struct}, sin
    indirección. Cero allocs para vértices de grado bajo (la mayoría
    en grafos scale-free). Cabe en una cache line.
    \item \textbf{Vec}: cuando el grado pasa 8, se promueve a
    \texttt{Vec<StreamEdge>}. Un alloc, acceso secuencial óptimo.
    El matcher hot-path lee el slice contiguo con
    \texttt{partition\_point} para acotar ventanas temporales en
    \(O(\log d)\) por frame.
    \item \textbf{Mmap}: cuando el budget total cruza el umbral
    (default 1\,M aristas, configurable vía
    \texttt{with\_memory\_budget\_edges}), todas las
    \texttt{NeighborList} \emph{Inline}/\emph{Vec} se vuelcan a un
    archivo memory-mapped sorted-by-timestamp y los slots pasan a
    \texttt{Mmap \{ handle, byte\_offset, len \}}. El \texttt{Arc<MmapHandle>}
    mantiene el archivo mapeado mientras quede al menos una
    referencia. Lecturas siguen siendo zero-copy: \texttt{as\_slice()}
    sobre \texttt{Mmap} proyecta el rango como \texttt{\&[StreamEdge]}.
\end{itemize}

\subsection{Por qué este shape y no la antigua \texttt{Btree}}

Una iteración previa del diseño tenía una variante \texttt{Btree}
para hubs (>128 aristas) que daba \(O(\log d)\) en queries por
ventana temporal. Se eliminó en \emph{stage 2c part 2} a favor de
\texttt{Vec} contiguo con \texttt{partition\_point}: el matcher
sostiene un slice borrowed sobre toda la duración del frame del DFS
con prefetch lookahead, y un \texttt{BTreeMap} no puede exponer un
slice contiguo. La pérdida algorítmica (\(O(d)\) lineal vs
\(O(\log d)\)) en queries por ventana se materializa solo en
consumers de \texttt{iter\_window} que ningún path de producción
usaba; el matcher hot-path tiene su propio \texttt{partition\_point}
sobre el slice \emph{Vec}/\emph{Mmap} que sigue dando \(O(\log d)\).

En grado bajo, \texttt{Inline} elimina el header del Vec (24~B)
+ heap alloc + indirección — para 1~M vértices con grado 1, eso son
24~MB salvados.

\section{Promotion + auto-spill}

Las transiciones entre variantes son atómicas; cada \texttt{append\_edge}
recorre tres puntos:

\begin{lstlisting}[language=Rust]
fn append_edge(&mut self, edge: StreamEdge) {
    match self {
        NeighborList::Inline(sv) if sv.len() < 8 => {
            sv.push(edge);
        }
        NeighborList::Inline(sv) => {
            // promote to Vec
            let mut v: Vec<StreamEdge> = sv.iter().copied().collect();
            v.push(edge);
            *self = NeighborList::Vec(v);
        }
        NeighborList::Vec(v) => {
            v.push(edge);
        }
        NeighborList::Mmap { .. } => {
            // post-spill write: copy out, append, swap
            let mut v: Vec<StreamEdge> = self.as_slice().to_vec();
            v.push(edge);
            *self = NeighborList::Vec(v);
        }
    }
}
\end{lstlisting}

\noindent
\emph{Inline \(\to\) Vec}: O(8) memcpy una vez por vértice cuando
cruza grado 8. Amortizado O(1).

\emph{Vec \(\to\) Mmap}: no es per-vertex sino un evento
\emph{global} disparado por el budget tracker
\texttt{edges\_since\_last\_spill}. El backend hace \texttt{auto\_spill()}
que sortea cada per-vertex Vec por timestamp, escribe un archivo
round-numbered, y swap-ea cada slot a \texttt{Mmap}. Costo: O(E)
bytes de I/O secuencial, una vez por round.

\emph{Mmap \(\to\) Vec}: O(d) memcpy out of mmap al primer write
post-spill sobre ese vértice. Producción no escribe post-spill
en el caso normal (toda la ingesta termina antes de que el
analista corra hunts), así que es cold path.

\section{Wait-free reads: \texttt{neighbors\_snapshot}}

La pieza interesante: el matcher quiere leer
\texttt{NeighborList} mientras el ingestor escribe. Las opciones
clásicas:

\begin{itemize}
    \item \texttt{Mutex}: serializa todo. Lento bajo lectura
    intensiva.
    \item \texttt{RwLock}: múltiples readers, pero el writer
    bloquea a todos los readers. Si la promotion es lenta, el
    matcher se traba.
    \item \emph{Wait-free}: los readers nunca bloquean. Esquema
    típico: el writer construye el nuevo estado en un buffer
    separado y publica el puntero atómicamente; los readers ven el
    puntero viejo o el nuevo, nunca un estado intermedio.
\end{itemize}

\noindent
El streaming backend implementa la tercera opción usando
\texttt{ArcSwap<NeighborList>}: cada vértice tiene un puntero
atómico a su \texttt{Arc<NeighborList>}. \texttt{neighbors\_snapshot}
hace un \texttt{load} atómico del puntero y devuelve el
\texttt{Arc} clonado:

\begin{lstlisting}[language=Rust]
pub fn neighbors_snapshot(&self, sid: StrId) -> Arc<NeighborList> {
    self.per_vertex[sid.0 as usize].load_full()
}
\end{lstlisting}

\noindent
El matcher ahora tiene su propia copia (refcount-shared) del
\texttt{NeighborList} que existía en el momento del
\texttt{load\_full()}. El writer puede seguir publicando nuevas
versiones; el matcher ve la suya \emph{consistente} hasta que
descarte el \texttt{Arc}.

\subsection{Costo}

\begin{itemize}
    \item Reader: 1 load atómico + 1 \texttt{Arc::clone}
    (\texttt{fetch\_add}) = $\sim 10$~ns total. Sin bloqueo.
    \item Writer: construir nuevo estado (puede ser $O(\deg)$ por
    promotion), 1 store atómico para publicar. El reader que vino
    \emph{antes} del store ve la versión vieja; el reader que vino
    \emph{después}, la nueva. No hay raza visible.
\end{itemize}

\subsection{Memory reclamation}

El reto de wait-free es \emph{cuándo liberar} el estado viejo. Si
un reader todavía tiene un \texttt{Arc} a la versión vieja, el
writer no puede liberarla. \texttt{ArcSwap} resuelve esto con
\emph{deferred drop}: el viejo \texttt{Arc} se desreferencia al
liberar el \texttt{ArcSwap}; el reader es responsable de soltar su
clone cuando termine. El refcount atómico hace el resto.

\section{¿Cómo se usa hoy?}

Post stage 2c part 2 el streaming backend es \emph{el} edge store.
Cinco \emph{ingest paths} comparten un solo helper
\texttt{push\_to\_streaming}: \texttt{add\_relation},
\texttt{add\_relations\_batch}, \texttt{replay\_relations},
\texttt{compact\_before}'s rebuild,
\texttt{remove\_entities\_and\_relations\_by\_dataset}'s rebuild.

Lectores production:

\begin{itemize}
    \item \textbf{Matcher hot path} (\texttt{dfs\_match\_smart},
    \texttt{dfs\_match\_iterative}, \texttt{edge\_range\_sid},
    AMAC pair): \texttt{neighbors\_arc(sid) + read guard +
    as\_slice()}. Slice borrowed sobre toda la duración del frame
    del DFS, con \texttt{partition\_point} para acotar ventanas
    temporales.
    \item \textbf{Index builds} (NLF, khop\_reach, lftj,
    yannakakis): \texttt{neighbors\_arc} per-vertex preserva cache
    locality sobre los acumuladores indexados por \texttt{v\_idx}.
    \item \textbf{Aggregaciones admin} (anomaly scorer warm-up,
    temporal heatmap, timeline data, schema cache fingerprint):
    \texttt{streaming.for\_each\_edge(closure)} con un solo
    \emph{outer guard} para todo el sweep.
    \item \textbf{Live-tail UI} (\texttt{cmd\_get\_recent\_edges})
    y telemetría (\texttt{cmd\_get\_ingest\_stats}): consumers
    delta-aware originales que sobrevivieron la migración.
\end{itemize}

\begin{enelcodigo}
\filepath{core/graph-engine/src/streaming.rs} expone el módulo
completo: \texttt{InMemoryStreamingBackend} con
\texttt{append\_edge}, \texttt{finalize}, \texttt{auto\_spill},
\texttt{neighbors\_arc}, \texttt{for\_each\_edge},
\texttt{for\_each\_vertex}, \texttt{consolidate}, \texttt{clear}.
El antiguo \texttt{spill.rs} con \texttt{SpillableEdgeStore} se
eliminó (stage 2c part 2 cleanup).
\end{enelcodigo}

\section{Resumen}

\begin{itemize}
    \item Streaming backend = el \emph{único} edge store, per-vertex
    con \texttt{NeighborList} \emph{degree-adaptive}:
    \texttt{Inline} ($\leq 8$), \texttt{Vec} (resto in-RAM),
    \texttt{Mmap} (post-spill).
    \item Auto-spill incremental dispara cuando el budget global
    (default $10^6$ aristas) cruza el umbral; sortea cada slice
    \emph{durante el write}, así el archivo es matcher-ready.
    \item \emph{Post-spill} writes promueven \texttt{Mmap} \(\to\)
    \texttt{Vec} sólo en el vértice tocado: cold path, una memcpy
    out of mmap.
    \item Lecturas son zero-copy via \texttt{neighbors\_arc + read
    guard + as\_slice}: el matcher hot-path sostiene el slice
    sobre toda la duración del frame, igual que el legacy
    \texttt{\&[CompactRelation]}.
    \item \texttt{StreamEdge} es layout-compatible con
    \texttt{CompactRelation} (\texttt{\#[repr(C)]}, mismo orden),
    así que los $\sim$50 readers del antiguo dual-write iteran el
    nuevo slice con field-access idéntico.
\end{itemize}

\begin{ejercicios}
\begin{enumerate}
    \item \starred{} Para un grafo de 1~M vértices distribuidos
    como power-law (\(\alpha = 2\)), estimá: cuántos vértices van a
    cada variante (\texttt{Inline}, \texttt{Vec}, \texttt{Btree}) y
    el ahorro de memoria vs usar \texttt{Vec} para todos.

    \item ¿Por qué la promotion es $O(\deg)$ \emph{una vez} por
    umbral cruzado y no $O(\deg)$ por insert? Mostrá la cuenta
    amortizada.

    \item \starred{} Diseñá un escenario donde \texttt{ArcSwap}
    falla (algo no es wait-free como prometido). (Pista: pensá en
    backpressure si los readers no descartan los Arcs viejos.)

    \item Compará: 1 millón de \texttt{neighbors\_snapshot} con
    \texttt{ArcSwap} vs con \texttt{RwLock<NeighborList>}. Tiempo y
    contention. ¿Cuándo gana cada uno?

    \item Si el motor migra a streaming como source of truth
    (stage 2c), el matcher hot loop pasa de iterar un slice
    contiguo a iterar una \texttt{NeighborList} variant. ¿Qué
    riesgo de regresión hay? ¿Para qué tipo de hipótesis?
\end{enumerate}
\end{ejercicios}
